\documentclass[11pt,a4paper]{article}

\usepackage[margin=25mm]{geometry}
\usepackage[T1]{fontenc}
\usepackage{lmodern}
\usepackage{microtype}
\usepackage{amsmath,amssymb,bm}
\usepackage{booktabs}
\usepackage{longtable}
\usepackage{array}
\usepackage{tabularx}
\usepackage{multirow}
\usepackage{enumitem}
\usepackage{xcolor}
\usepackage{listings}
\usepackage[colorlinks=true,linkcolor=blue!70!black,citecolor=blue!70!black,urlcolor=blue!70!black]{hyperref}
\usepackage{natbib}
\bibliographystyle{plainnat}

% Code style
\definecolor{codebg}{RGB}{248,248,248}
\definecolor{codeframe}{RGB}{200,200,200}
\definecolor{codecomment}{RGB}{0,120,0}
\definecolor{codekeyword}{RGB}{0,0,160}
\definecolor{codestring}{RGB}{160,20,20}

\lstdefinestyle{Rstyle}{%
  language=R,
  backgroundcolor=\color{codebg},
  frame=single,
  rulecolor=\color{codeframe},
  basicstyle=\ttfamily\small,
  keywordstyle=\color{codekeyword}\bfseries,
  commentstyle=\color{codecomment}\itshape,
  stringstyle=\color{codestring},
  breaklines=true,
  showstringspaces=false,
  tabsize=2,
  numbers=none,
  morekeywords={bdeb_data,bdeb_model,bdeb_fit,bdeb_diagnose,bdeb_ppc,
    bdeb_summary,bdeb_derived,bdeb_predict,bdeb_tox,bdeb_ec50,
    prior_lognormal,prior_normal,prior_beta,prior_halfnormal,
    prior_halfcauchy,prior_exponential,prior_default,
    obs_normal,obs_lognormal,obs_student_t,obs_poisson,obs_negbinom,
    arrhenius,deb_fluxes,repro_to_intervals,plot_dose_response,
    library,data,head,plot,summary}
}
\lstset{style=Rstyle}

\lstdefinestyle{Stanstyle}{%
  backgroundcolor=\color{codebg},
  frame=single,
  rulecolor=\color{codeframe},
  basicstyle=\ttfamily\small,
  keywordstyle=\color{codekeyword}\bfseries,
  commentstyle=\color{codecomment}\itshape,
  breaklines=true,
  showstringspaces=false,
  tabsize=2,
  numbers=none,
  morekeywords={functions,data,parameters,transformed,model,generated,
    quantities,real,int,vector,array,lower,upper,ode_bdf,normal,
    lognormal,beta,neg_binomial_2,poisson,fmax,pow,exp,log}
}

% Custom commands
\newcommand{\pkg}[1]{\textbf{#1}}
\newcommand{\fn}[1]{\texttt{#1()}}
\newcommand{\code}[1]{\texttt{#1}}
\newcommand{\R}{\textsf{R}}
\newcommand{\Stan}{\textsf{Stan}}
\newcommand{\BDEB}{\textit{Bayesian DEB}}

\emergencystretch=2em

% ============================================================
\title{\pkg{BayesianDEB}: Bayesian Dynamic Energy Budget Modelling in \R\\[1ex]
       \large Package Manual and Tutorial (Version 0.1.0)}
\author{Branimir K.\ Hackenberger}
\date{April 2026}

\begin{document}
\maketitle

\begin{abstract}
\pkg{BayesianDEB} is an \R{} package that provides a complete Bayesian
framework for Dynamic Energy Budget (DEB) modelling.  It wraps pre-written
\Stan{} models with a declarative \R{} interface for data preparation,
prior specification, MCMC fitting via \pkg{cmdstanr}, convergence
diagnostics, posterior predictive checks, derived quantity estimation, and
publication-quality visualisation.  Four model types are included:
individual-level growth, growth with reproduction, hierarchical
multi-individual analysis with partial pooling, and DEBtox
(toxicokinetic-toxicodynamic) models for ecotoxicology.  This document
provides a comprehensive reference for all exported functions, bundled Stan
models, example datasets, and a step-by-step tutorial covering the full
\BDEB{} workflow.
\end{abstract}

\tableofcontents
\newpage

% ============================================================
\section{Introduction}
\label{sec:introduction}

Dynamic Energy Budget (DEB) theory \citep{Kooijman2010} provides a
mechanistic, thermodynamically consistent description of how organisms
acquire and utilise energy for maintenance, growth, development, and
reproduction.  Classical DEB calibration relies on deterministic
optimisation---finding a single best-fit parameter vector---which yields no
formal uncertainty quantification and cannot naturally accommodate prior
biological knowledge or hierarchical data structures.

\pkg{BayesianDEB} addresses these limitations by embedding the DEB
differential equation system within a Bayesian state-space framework.
Parameters are treated as random variables with prior distributions; the
likelihood is constructed from observation models linking latent DEB states
to measured quantities; and the posterior is explored via Hamiltonian Monte
Carlo (HMC) using \Stan{} \citep{Carpenter2017}.

The package workflow follows the iterative Bayesian modelling cycle
advocated by \citet{Gelman2020}:
\begin{enumerate}[nosep]
  \item Prepare and validate data.
  \item Specify the DEB process model, priors, and observation model.
  \item Fit the model via MCMC.
  \item Diagnose convergence.
  \item Perform posterior predictive checks.
  \item Compute derived quantities and predictions.
  \item Revise the model as needed.
\end{enumerate}

% ============================================================
\section{Installation}
\label{sec:installation}

\pkg{BayesianDEB} requires \pkg{cmdstanr} and a working CmdStan
installation for model fitting.  Data preparation, prior specification,
and utility functions work without \Stan{}.

\begin{lstlisting}
# 1. Install cmdstanr from Stan r-universe
install.packages("cmdstanr",
  repos = c("https://stan-dev.r-universe.dev",
            getOption("repos")))

# 2. Install CmdStan (one-time, ~2 min)
cmdstanr::install_cmdstan()

# 3. Install BayesianDEB
install.packages("BayesianDEB")
# or from source:
# remotes::install_github("bhackenberger/BayesianDEB")

# 4. Load
library(BayesianDEB)
\end{lstlisting}

% ============================================================
\section{DEB Theory Recap}
\label{sec:deb-recap}

The standard DEB model tracks four state variables over time:
\begin{itemize}[nosep]
  \item $E$ --- reserve energy (J),
  \item $V$ --- structural volume (cm$^3$), with structural length $L = V^{1/3}$,
  \item $E_H$ --- maturity (J),
  \item $E_R$ --- reproduction buffer (J).
\end{itemize}

Energy flows are governed by the $\kappa$-rule: a fraction $\kappa$ of
mobilised reserve goes to somatic maintenance and growth; the remainder
goes to maturity maintenance and reproduction.  The key fluxes are:
\begin{align}
  \dot{p}_A &= f\,\{p_{Am}\}\,L^2  &&\text{(assimilation)}, \label{eq:pA}\\
  \dot{p}_C &= \frac{E\,v\,L}{E + [E_G]\,V}  &&\text{(mobilisation)}, \label{eq:pC}\\
  \dot{p}_M &= [p_M]\,V  &&\text{(somatic maintenance)}, \label{eq:pM}\\
  \frac{\mathrm{d}E}{\mathrm{d}t} &= \dot{p}_A - \dot{p}_C,  \label{eq:dE}\\
  \frac{\mathrm{d}V}{\mathrm{d}t} &= \frac{\kappa\,\dot{p}_C - \dot{p}_M}{[E_G]},  \label{eq:dV}\\
  \frac{\mathrm{d}E_R}{\mathrm{d}t} &= (1-\kappa)\,\dot{p}_C - \dot{p}_J.  \label{eq:dR}
\end{align}

Table~\ref{tab:params} lists the standard DEB parameters used throughout
this package.

\begin{table}[htbp]
\centering
\caption{Standard DEB parameters in \pkg{BayesianDEB}.}
\label{tab:params}
\begin{tabular}{llll}
\toprule
\textbf{Symbol} & \textbf{R name} & \textbf{Units} & \textbf{Description} \\
\midrule
$\{p_{Am}\}$ & \code{p\_Am} & J\,d$^{-1}$\,cm$^{-2}$ & Max.\ surface-area-specific assimilation \\
$[p_M]$      & \code{p\_M}  & J\,d$^{-1}$\,cm$^{-3}$ & Volume-specific somatic maintenance \\
$\kappa$     & \code{kappa} & -- & Allocation fraction to soma \\
$v$          & \code{v}     & cm\,d$^{-1}$ & Energy conductance \\
$[E_G]$      & \code{E\_G}  & J\,cm$^{-3}$ & Specific cost of structure \\
$k_J$        & \code{k\_J}  & d$^{-1}$ & Maturity maintenance rate \\
$E_H^p$      & \code{E\_Hp} & J & Maturity at puberty \\
$f$          & \code{f\_food} & -- & Scaled functional response (0--1) \\
$E_0$        & \code{E0}    & J\,cm$^{-3}$ & Initial reserve density \\
$L_0$        & \code{L0}    & cm & Initial structural length \\
$\sigma_L$   & \code{sigma\_L} & cm & Observation error SD (length) \\
\bottomrule
\end{tabular}
\end{table}

% ============================================================
\section{Package Architecture}
\label{sec:architecture}

\pkg{BayesianDEB} consists of three layers:

\begin{description}[style=nextline]
  \item[Stan models] Pre-written \code{.stan} files in \code{inst/stan/}
    implementing the DEB ODE system and Bayesian inference.  Users never
    need to edit Stan code directly.
  \item[R interface] Exported functions for data preparation, model
    specification, fitting, diagnostics, and visualisation.  All functions
    use S3 classes (\code{bdeb\_data}, \code{bdeb\_model}, \code{bdeb\_fit}).
  \item[Prior system] A family of \code{prior\_*()} constructor functions
    that encode distribution choices as lightweight objects passed to
    \fn{bdeb\_model}.
\end{description}

\subsection{Bundled Stan Models}

\begin{table}[htbp]
\centering
\caption{Stan models bundled with \pkg{BayesianDEB}.}
\label{tab:stan-models}
\begin{tabular}{lccl}
\toprule
\textbf{Stan file} & \textbf{States} & \textbf{R type} & \textbf{Use case} \\
\midrule
\code{bdeb\_individual\_growth} & $E, V$ & \code{"individual"} & Single organism growth \\
\code{bdeb\_growth\_repro} & $E, V, R$ & \code{"growth\_repro"} & Growth + reproduction \\
\code{bdeb\_hierarchical\_growth} & $E, V$ + RE & \code{"hierarchical"} & Multi-individual, partial pooling \\
\code{bdeb\_debtox} & $E, V, R, D_w$ & \code{"debtox"} & Ecotoxicology (TKTD) \\
\bottomrule
\end{tabular}
\end{table}

All Stan models use the BDF (backward differentiation formula) stiff ODE
solver (\code{ode\_bdf}) with tolerances $10^{-6}$ and maximum 10\,000
steps, appropriate for the time-scale separation typical of DEB reserve
and structural dynamics.

% ============================================================
\section{Complete Function Reference}
\label{sec:functions}

% ---- 5.1 Data preparation ----
\subsection{Data Preparation}

\subsubsection{\fn{bdeb\_data} --- Prepare Data for BDEB Models}

\begin{lstlisting}
bdeb_data(growth = NULL, reproduction = NULL, survival = NULL,
          concentration = NULL, f_food = 1.0)
\end{lstlisting}

Converts long-format data frames into the structured format required by
the Stan models.  Validates column names, sorts by individual and time,
and constructs the internal \code{bdeb\_data} S3 object.

\paragraph{Arguments.}
\begin{description}[nosep,leftmargin=3cm,style=nextline]
  \item[\code{growth}] Data frame with columns \code{id}, \code{time}, \code{length}.
  \item[\code{reproduction}] Data frame with columns \code{id}, \code{t\_start}, \code{t\_end}, \code{count}.
  \item[\code{survival}] Data frame with columns \code{id}, \code{time}, \code{event} (1~=~death, 0~=~censored).
  \item[\code{concentration}] Named numeric vector or data frame mapping IDs to external toxicant concentrations (for DEBtox).
  \item[\code{f\_food}] Scaled functional response, $f \in [0,1]$.  Default 1 (ad libitum).
\end{description}

\paragraph{Returns.} A \code{bdeb\_data} object (list with class \code{"bdeb\_data"}).

\paragraph{Example.}
\begin{lstlisting}
df <- data.frame(
  id = rep(1, 10),
  time = seq(0, 45, by = 5),
  length = c(0.10, 0.15, 0.22, 0.30, 0.38,
             0.44, 0.49, 0.52, 0.54, 0.55)
)
dat <- bdeb_data(growth = df)
dat
\end{lstlisting}

\subsubsection{\fn{repro\_to\_intervals} --- Convert Cumulative Reproduction to Intervals}

\begin{lstlisting}
repro_to_intervals(df)
\end{lstlisting}

Converts a data frame of cumulative offspring counts to interval counts
suitable for the \code{reproduction} argument of \fn{bdeb\_data}.

\paragraph{Arguments.}
\begin{description}[nosep,leftmargin=3cm]
  \item[\code{df}] Data frame with columns \code{id}, \code{time}, \code{cumulative}.
\end{description}

\paragraph{Returns.} Data frame with columns \code{id}, \code{t\_start}, \code{t\_end}, \code{count}.

\paragraph{Example.}
\begin{lstlisting}
cumul <- data.frame(
  id = rep(1, 5),
  time = c(0, 7, 14, 21, 28),
  cumulative = c(0, 10, 30, 60, 100)
)
intervals <- repro_to_intervals(cumul)
intervals
#   id t_start t_end count
# 1  1       0     7    10
# 2  1       7    14    20
# 3  1      14    21    30
# 4  1      21    28    40
\end{lstlisting}

% ---- 5.2 Priors ----
\subsection{Prior Specification}
\label{sec:priors}

Prior distributions are specified via lightweight constructor functions
that return \code{bdeb\_prior} objects.  These are passed as a named list
to \fn{bdeb\_model}; any unspecified priors are filled from
\fn{prior\_default}.

\subsubsection{Prior Constructors}

\begin{longtable}{lll}
\toprule
\textbf{Function} & \textbf{Distribution} & \textbf{Arguments} \\
\midrule
\endhead
\fn{prior\_lognormal} & $\text{LogNormal}(\mu, \sigma)$ & \code{mu = 0, sigma = 1} \\
\fn{prior\_normal} & $\mathcal{N}(\mu, \sigma)$ & \code{mu = 0, sigma = 1} \\
\fn{prior\_beta} & $\text{Beta}(a, b)$ & \code{a = 2, b = 2} \\
\fn{prior\_halfnormal} & $\text{HalfNormal}(\sigma)$ & \code{sigma = 1} \\
\fn{prior\_halfcauchy} & $\text{HalfCauchy}(\sigma)$ & \code{sigma = 1} \\
\fn{prior\_exponential} & $\text{Exponential}(\lambda)$ & \code{rate = 1} \\
\bottomrule
\end{longtable}

All return a \code{bdeb\_prior} object (list with elements \code{family}
and the hyperparameters).

\paragraph{Usage recommendations.}
\begin{itemize}[nosep]
  \item Positive rate parameters ($\{p_{Am}\}$, $[p_M]$, $v$, $[E_G]$): \fn{prior\_lognormal}.
  \item Allocation fraction ($\kappa$): \fn{prior\_beta}.
  \item Observation error SD ($\sigma_L$): \fn{prior\_halfnormal}.
  \item Hierarchical variance components: \fn{prior\_exponential}.
\end{itemize}

\subsubsection{\fn{prior\_default} --- Weakly Informative Defaults}

\begin{lstlisting}
prior_default(type = c("individual", "growth_repro",
                        "hierarchical", "debtox"))
\end{lstlisting}

Returns a complete named list of weakly informative priors appropriate for
the specified model type.  Table~\ref{tab:default-priors} shows the
defaults for the \code{"individual"} type.

\begin{table}[htbp]
\centering
\caption{Default priors for individual growth model.}
\label{tab:default-priors}
\begin{tabular}{lll}
\toprule
\textbf{Parameter} & \textbf{Prior} & \textbf{Rationale} \\
\midrule
\code{p\_Am}    & $\text{LogNormal}(1.0,\; 1.0)$   & Covers range $\sim$0.4--20 \\
\code{p\_M}     & $\text{LogNormal}(-1.0,\; 1.0)$  & Covers range $\sim$0.04--2 \\
\code{kappa}    & $\text{Beta}(2,\; 2)$             & Favours interior values \\
\code{v}        & $\text{LogNormal}(-1.5,\; 1.0)$   & Covers range $\sim$0.02--1 \\
\code{E\_G}     & $\text{LogNormal}(6.0,\; 1.0)$    & Covers $\sim$100--3000 \\
\code{E0}       & $\text{LogNormal}(0.0,\; 1.0)$    & \\
\code{L0}       & $\text{LogNormal}(-2.0,\; 1.0)$   & Initial length $\sim$0.01--0.5 \\
\code{sigma\_L} & $\text{HalfNormal}(0.1)$          & Measurement error SD \\
\bottomrule
\end{tabular}
\end{table}

Additional priors are included for specific types:
\begin{itemize}[nosep]
  \item \code{"growth\_repro"} adds: \code{k\_J}, \code{E\_Hp}, \code{k\_R}, \code{phi\_R}.
  \item \code{"hierarchical"} adds: \code{mu\_log\_p\_Am} (Normal), \code{sigma\_log\_p\_Am} (Exponential).
  \item \code{"debtox"} adds: \code{k\_d}, \code{z\_w}, \code{b\_w}.
\end{itemize}

% ---- 5.3 Model specification ----
\subsection{Model Specification}

\subsubsection{\fn{bdeb\_model} --- Specify a BDEB Model}

\begin{lstlisting}
bdeb_model(data, type = c("individual", "growth_repro",
                           "hierarchical", "debtox"),
           priors = list(), observation = list(),
           temperature = NULL)
\end{lstlisting}

Creates a model specification combining data, DEB model type, observation
model, and prior distributions.  Does not fit the model.

\paragraph{Arguments.}
\begin{description}[nosep,leftmargin=3.5cm,style=nextline]
  \item[\code{data}] A \code{bdeb\_data} object.
  \item[\code{type}] Model type (see Table~\ref{tab:stan-models}).
  \item[\code{priors}] Named list of \code{bdeb\_prior} objects.  Missing entries filled from \fn{prior\_default}.
  \item[\code{observation}] Named list: \code{growth = obs\_normal()}, \code{reproduction = obs\_negbinom()}, etc.
  \item[\code{temperature}] Optional list with \code{T}, \code{T\_ref}, \code{T\_A} for Arrhenius correction.
\end{description}

\paragraph{Returns.} A \code{bdeb\_model} object.

\subsubsection{Observation Model Constructors}

\begin{longtable}{lll}
\toprule
\textbf{Function} & \textbf{Distribution} & \textbf{Default for} \\
\midrule
\fn{obs\_normal}    & $\mathcal{N}(\hat{L}, \sigma_L)$ & Growth (continuous) \\
\fn{obs\_lognormal} & $\text{LogNormal}(\log\hat{L}, \sigma)$ & Multiplicative error \\
\fn{obs\_student\_t} & $t_\nu(\hat{L}, \sigma_L)$ & Outlier-robust \\
\fn{obs\_poisson}   & $\text{Poisson}(\Delta R)$ & Count data \\
\fn{obs\_negbinom}  & $\text{NegBin}(\Delta R, \phi)$ & Reproduction (default) \\
\bottomrule
\end{longtable}

% ---- 5.4 Fitting ----
\subsection{Model Fitting}

\subsubsection{\fn{bdeb\_fit} --- Fit a BDEB Model}

\begin{lstlisting}
bdeb_fit(model, chains = 4, iter_warmup = 1000,
         iter_sampling = 1000, adapt_delta = 0.8,
         max_treedepth = 10, seed = NULL,
         parallel_chains = NULL, refresh = 200, ...)
\end{lstlisting}

Compiles the Stan model and runs MCMC sampling via \pkg{cmdstanr}.

\paragraph{Key arguments.}
\begin{description}[nosep,leftmargin=3.5cm]
  \item[\code{model}] A \code{bdeb\_model} object.
  \item[\code{adapt\_delta}] Target acceptance probability.  Increase toward 1.0 if divergences occur.
  \item[\code{max\_treedepth}] Maximum NUTS tree depth.  Increase if treedepth is saturated.
  \item[\code{seed}] Integer seed for reproducibility.
\end{description}

\paragraph{Returns.} A \code{bdeb\_fit} object containing the raw
\code{CmdStanMCMC} result (\code{fit\$fit}), the model specification
(\code{fit\$model}), and sampling metadata.

% ---- 5.5 Diagnostics ----
\subsection{Diagnostics and Posterior Analysis}

\subsubsection{\fn{bdeb\_diagnose} --- MCMC Convergence Diagnostics}

\begin{lstlisting}
bdeb_diagnose(fit, pars = NULL)
\end{lstlisting}

Reports: divergent transitions, treedepth saturation, E-BFMI, $\hat{R}$,
bulk and tail ESS for each parameter.  Flags problematic values
automatically.

\paragraph{Returns.} Invisibly, a list with \code{n\_divergent},
\code{n\_max\_treedepth}, \code{ebfmi}, and \code{summary} (a
\code{posterior::draws\_summary} data frame).

\subsubsection{\fn{bdeb\_summary} --- Posterior Summary Table}

\begin{lstlisting}
bdeb_summary(fit, pars = NULL, prob = 0.90)
\end{lstlisting}

Returns a tidy summary of posterior draws: mean, SD, median, credible
interval bounds, $\hat{R}$, and ESS.

\paragraph{Arguments.}
\begin{description}[nosep,leftmargin=2.5cm]
  \item[\code{pars}] Character vector of parameter names.  Default: all.
  \item[\code{prob}] Credible interval width.  Default 0.90 (5th--95th percentiles).
\end{description}

\subsubsection{\fn{bdeb\_derived} --- Derived Biological Quantities}

\begin{lstlisting}
bdeb_derived(fit, quantities = c("L_inf", "k_M",
             "growth_rate"), f = 1.0)
\end{lstlisting}

Computes derived quantities from posterior draws, propagating uncertainty:
\begin{itemize}[nosep]
  \item \code{"L\_inf"}: ultimate structural length, $L_\infty = f\,\kappa\,\{p_{Am}\} / [p_M]$.
  \item \code{"k\_M"}: somatic maintenance rate constant, $k_M = [p_M] / [E_G]$.
  \item \code{"growth\_rate"}: von Bertalanffy rate, $\dot{k} = \frac{v}{3}\,\frac{[p_M]}{\kappa\,[E_G]}$.
\end{itemize}

\paragraph{Returns.} A \code{posterior::draws\_df} with one column per
requested quantity and one row per posterior draw.

\subsubsection{\fn{bdeb\_ppc} --- Posterior Predictive Checks}

\begin{lstlisting}
bdeb_ppc(fit, type = c("growth", "reproduction", "all"))
\end{lstlisting}

Extracts replicated observations from the \code{generated quantities}
block of the Stan model.

\paragraph{Returns.} A \code{bdeb\_ppc} object with matrices of
replicated data (\code{L\_rep}, \code{R\_rep}) and the corresponding
observed data.

\subsubsection{\fn{bdeb\_predict} --- Posterior Predictions}

\begin{lstlisting}
bdeb_predict(fit, newdata = NULL, n_draws = 200)
\end{lstlisting}

Extracts posterior predicted trajectories.  Can also be called as
\code{predict(fit)}.

% ---- 5.6 Plotting ----
\subsection{Plotting}

\subsubsection{\code{plot(fit, type = ...)} --- Plot a BDEB Fit}

\begin{lstlisting}
plot(fit, type = c("trace", "posterior", "pairs",
                    "trajectory"),
     pars = NULL, n_draws = 100)
\end{lstlisting}

Available plot types:
\begin{description}[nosep,leftmargin=3cm]
  \item[\code{"trace"}] MCMC trace plots for convergence assessment.
  \item[\code{"posterior"}] Overlaid marginal posterior density curves.
  \item[\code{"pairs"}] Bivariate scatter plots showing posterior correlations.
  \item[\code{"trajectory"}] Posterior predicted growth curves overlaid on observed data.
\end{description}

All return \code{ggplot2} objects.

\subsubsection{\code{plot(ppc)} --- Plot Posterior Predictive Checks}

\begin{lstlisting}
ppc <- bdeb_ppc(fit)
plot(ppc, n_draws = 50)
\end{lstlisting}

Overlays replicated trajectories (grey) on observed data points (red).

% ---- 5.7 DEBtox ----
\subsection{DEBtox (Ecotoxicology)}

\subsubsection{\fn{bdeb\_tox} --- DEBtox Model Specification}

\begin{lstlisting}
bdeb_tox(data, stress = c("assimilation", "maintenance",
         "growth_cost"), priors = list(), ...)
\end{lstlisting}

Convenience wrapper for \fn{bdeb\_model} with \code{type = "debtox"}.
Currently implements stress on assimilation (reduces $\dot{p}_A$ by a
factor $\max(1 - s, 0)$ where $s = b_w \max(D_w, 0)$).

The internal damage $D_w$ follows:
\begin{equation}
\frac{\mathrm{d}D_w}{\mathrm{d}t} = k_d\bigl(\max(C_w - z_w, 0) - D_w\bigr),
\end{equation}
where $k_d$ is the damage recovery rate, $z_w$ is the no-effect
concentration (NEC), and $C_w$ is external concentration.

\subsubsection{\fn{bdeb\_ec50} --- Extract EC50 and NEC}

\begin{lstlisting}
bdeb_ec50(fit, prob = 0.90)
\end{lstlisting}

Extracts the posterior distribution of the $\text{EC}_{50}$ computed
analytically in Stan as:
\begin{equation}
\text{EC}_{50} = z_w + \frac{0.5}{b_w}.
\end{equation}

\paragraph{Returns.} A list with \code{draws} (posterior vector),
\code{summary} (data frame with mean, median, SD, credible interval), and
\code{NEC} (summary of the no-effect concentration).

\subsubsection{\fn{plot\_dose\_response} --- Dose-Response Curve}

\begin{lstlisting}
plot_dose_response(fit, endpoint = "growth", n_draws = 100)
\end{lstlisting}

Plots the stress function $\max(1 - b_w\max(C - z_w, 0),\; 0)$ over a
continuous concentration range with posterior uncertainty bands, overlaying
observed endpoint values.

% ---- 5.8 Utilities ----
\subsection{Utility Functions}

\subsubsection{\fn{arrhenius} --- Temperature Correction}

\begin{lstlisting}
arrhenius(T, T_ref = 293.15, T_A = 8000)
\end{lstlisting}

Returns the Arrhenius correction factor:
\begin{equation}
c_T = \exp\!\left(\frac{T_A}{T_{\text{ref}}} - \frac{T_A}{T}\right).
\end{equation}

\paragraph{Example.}
\begin{lstlisting}
# At reference temperature: factor = 1
arrhenius(293.15)
# [1] 1

# At 25C (5 degrees warmer): factor > 1
arrhenius(298.15)
# [1] 1.741
\end{lstlisting}

\subsubsection{\fn{deb\_fluxes} --- Compute DEB Energy Fluxes}

\begin{lstlisting}
deb_fluxes(E, V, f, p_Am, p_M, kappa, v, E_G,
           k_J = 0, E_Hp = 0)
\end{lstlisting}

Given current state $(E, V)$ and parameters, returns a named list of all
standard DEB fluxes: \code{p\_A}, \code{p\_C}, \code{p\_M}, \code{p\_G},
\code{p\_J}, \code{p\_R}, plus structural length \code{L} and scaled
reserve density~\code{e}.

\paragraph{Example.}
\begin{lstlisting}
fl <- deb_fluxes(E = 10, V = 0.5, f = 1.0,
                 p_Am = 5, p_M = 0.5, kappa = 0.75,
                 v = 0.2, E_G = 400)
fl$p_A   # assimilation flux
fl$L     # structural length = V^(1/3)
\end{lstlisting}


% ============================================================
\section{Bundled Datasets}
\label{sec:datasets}

\subsection{\code{eisenia\_growth}}

Simulated growth data for 21 individuals of the earthworm \textit{Eisenia
fetida} measured weekly over 12 weeks (days 0, 7, 14, \ldots, 84).

\begin{description}[nosep,leftmargin=3cm]
  \item[Rows] 273 (21 individuals $\times$ 13 time points).
  \item[Columns] \code{id} (integer 1--21), \code{time} (days), \code{length} (cm).
  \item[True parameters] $\{p_{Am}\} = 5.0$, $[p_M] = 0.5$, $\kappa = 0.75$, $v = 0.2$, $[E_G] = 400$; individual variation in $\{p_{Am}\}$ (CV $\approx$ 10\%) and $L_0$ (CV $\approx$ 5\%); $\sigma_L = 0.015$\,cm.
\end{description}

\subsection{\code{folsomia\_repro}}

Simulated reproduction data for the springtail \textit{Folsomia candida}
exposed to 5 cadmium concentrations (0, 10, 50, 100, 500\,mg/kg) with
6 replicates per concentration.  Standard 28-day test.

\begin{description}[nosep,leftmargin=3cm]
  \item[Rows] 30 (5 concentrations $\times$ 6 replicates).
  \item[Columns] \code{id}, \code{concentration}, \code{t\_start}, \code{t\_end}, \code{count}.
  \item[True parameters] NEC = 30\,mg/kg, $b_w = 0.005$, base offspring = 120.
\end{description}

\subsection{\code{debtox\_growth}}

Simulated growth under toxicant exposure: 4 concentration groups (0, 20,
80, 200), 10 individuals each, weekly measurements over 6 weeks.

\begin{description}[nosep,leftmargin=3cm]
  \item[Rows] 280 (40 individuals $\times$ 7 time points).
  \item[Columns] \code{id}, \code{concentration}, \code{time}, \code{length}.
  \item[True parameters] NEC = 15, $b_w = 0.003$, stress on assimilation.
\end{description}


% ============================================================
\section{Tutorial: Complete Workflow}
\label{sec:tutorial}

This section walks through all four model types using the bundled
datasets.

\subsection{Tutorial 1: Individual Growth}
\label{sec:tut-individual}

\begin{lstlisting}
library(BayesianDEB)
data(eisenia_growth)

# --- Step 1: Prepare data (single individual) ---
df1 <- eisenia_growth[eisenia_growth$id == 1, ]
dat <- bdeb_data(growth = df1, f_food = 1.0)
dat

# --- Step 2: Specify model ---
mod <- bdeb_model(dat, type = "individual",
  priors = list(
    p_Am    = prior_lognormal(mu = 1.5, sigma = 0.5),
    p_M     = prior_lognormal(mu = -1.0, sigma = 0.5),
    kappa   = prior_beta(a = 3, b = 2),
    sigma_L = prior_halfnormal(sigma = 0.05)
  ))
mod  # prints specification including all priors

# --- Step 3: Fit via MCMC ---
fit <- bdeb_fit(mod, chains = 4, iter_sampling = 2000,
                adapt_delta = 0.9, seed = 42)
fit  # prints quick diagnostic summary

# --- Step 4: Convergence diagnostics ---
bdeb_diagnose(fit)
# Check:
#   - No divergent transitions
#   - All R-hat < 1.01
#   - Bulk ESS > 400 for all parameters

# Trace plots
plot(fit, type = "trace")

# Posterior correlations
plot(fit, type = "pairs",
     pars = c("p_Am", "p_M", "kappa"))

# --- Step 5: Posterior summary ---
bdeb_summary(fit, prob = 0.95)

# --- Step 6: Derived quantities ---
bdeb_derived(fit, quantities = c("L_inf", "growth_rate"))

# --- Step 7: Posterior predictive check ---
ppc <- bdeb_ppc(fit, type = "growth")
plot(ppc, n_draws = 100)

# --- Step 8: Trajectory plot ---
plot(fit, type = "trajectory", n_draws = 200)
\end{lstlisting}

\subsection{Tutorial 2: Hierarchical Multi-Individual}
\label{sec:tut-hierarchical}

\begin{lstlisting}
data(eisenia_growth)

# All 21 individuals
dat_all <- bdeb_data(growth = eisenia_growth, f_food = 1.0)
dat_all
# => 21 individuals, 273 observations

# Hierarchical model: individual variation in p_Am
# Non-centred parameterisation (automatic in the Stan model)
mod_h <- bdeb_model(dat_all, type = "hierarchical",
  priors = list(
    mu_log_p_Am    = prior_normal(mu = 1.5, sigma = 0.5),
    sigma_log_p_Am = prior_exponential(rate = 2),
    kappa          = prior_beta(a = 3, b = 2)
  ))

fit_h <- bdeb_fit(mod_h, chains = 4, iter_sampling = 2000,
                  adapt_delta = 0.95, seed = 123)

# Population-level summary
bdeb_summary(fit_h,
  pars = c("mu_log_p_Am", "sigma_log_p_Am",
           "p_M", "kappa", "v", "E_G"))

# Individual assimilation rates
bdeb_summary(fit_h,
  pars = paste0("p_Am_ind[", 1:5, "]"))

# Prediction for a new individual
bdeb_summary(fit_h, pars = "p_Am_new")

# Trace plots for hierarchical parameters
plot(fit_h, type = "trace",
  pars = c("mu_log_p_Am", "sigma_log_p_Am"))
\end{lstlisting}

\subsection{Tutorial 3: Growth + Reproduction}
\label{sec:tut-growth-repro}

\begin{lstlisting}
# Simulated growth + reproduction data
growth_df <- data.frame(
  id = 1, time = seq(0, 56, by = 7),
  length = c(0.10, 0.18, 0.28, 0.37, 0.44,
             0.49, 0.52, 0.54, 0.55)
)

repro_df <- data.frame(
  id = 1,
  t_start = c(28, 35, 42, 49),
  t_end   = c(35, 42, 49, 56),
  count   = c(5, 12, 18, 15)
)

dat_gr <- bdeb_data(growth = growth_df,
                    reproduction = repro_df,
                    f_food = 1.0)

mod_gr <- bdeb_model(dat_gr, type = "growth_repro",
  priors = list(
    kappa = prior_beta(a = 3, b = 2),
    k_R   = prior_lognormal(mu = -2, sigma = 1)
  ))

fit_gr <- bdeb_fit(mod_gr, chains = 4,
                   adapt_delta = 0.95, seed = 99)

bdeb_diagnose(fit_gr)
bdeb_summary(fit_gr)

# Growth PPC
plot(bdeb_ppc(fit_gr, type = "growth"))
\end{lstlisting}

\subsection{Tutorial 4: DEBtox Ecotoxicology}
\label{sec:tut-debtox}

\begin{lstlisting}
data(debtox_growth)

# Concentration mapping: group means
conc_map <- c(0, 20, 80, 200)
names(conc_map) <- unique(debtox_growth$concentration)

dat_tox <- bdeb_data(
  growth = debtox_growth,
  concentration = conc_map,
  f_food = 1.0
)

# DEBtox model: stress on assimilation
mod_tox <- bdeb_tox(dat_tox, stress = "assimilation",
  priors = list(
    z_w = prior_lognormal(mu = 2.5, sigma = 1.0),
    b_w = prior_lognormal(mu = -5.0, sigma = 2.0)
  ))

fit_tox <- bdeb_fit(mod_tox, chains = 4,
                    adapt_delta = 0.95, seed = 77)

# Diagnostics
bdeb_diagnose(fit_tox)

# EC50 and NEC with 95% credible intervals
bdeb_ec50(fit_tox, prob = 0.95)
# Output:
#   parameter  mean median    sd lower upper
#       EC50  ...   ...    ...   ...   ...
#        NEC  ...   ...    ...   ...   ...

# Dose-response plot with posterior uncertainty
plot_dose_response(fit_tox, n_draws = 200)

# Parameter summary
bdeb_summary(fit_tox,
  pars = c("p_Am", "p_M", "kappa", "k_d", "z_w", "b_w"))
\end{lstlisting}


% ============================================================
\section{Stan Model Details}
\label{sec:stan-details}

This section describes the internal structure of the bundled Stan models
for users who wish to understand or extend them.

\subsection{DEB ODE Right-Hand Side}

All models share the same ODE function implementing Equations
\eqref{eq:pA}--\eqref{eq:dV}.  The function signature is:

\begin{lstlisting}[style=Stanstyle]
vector deb_growth_ode(real t, vector x,
  array[] real theta, array[] real env)
\end{lstlisting}

where \code{x = [E, V]'}, \code{theta = \{p\_Am, p\_M, kappa, v, E\_G\}},
and \code{env = \{f\}}.

A stabilisation term \code{1e-12} is added to the denominator of the
mobilisation flux to prevent division by zero:
\begin{equation}
\dot{p}_C = \frac{E\,v\,L}{E + [E_G]\,V + 10^{-12}}.
\end{equation}

\subsection{Prior Hyperparameters as Data}

Prior hyperparameters are passed from \R{} as Stan \code{data} variables
(e.g., \code{prior\_p\_Am\_mu}, \code{prior\_p\_Am\_sd}).  This allows
users to change priors without recompiling the Stan model.

\subsection{Hierarchical Non-Centred Parameterisation}

The hierarchical model uses the standard non-centred parameterisation to
avoid funnel geometries:
\begin{align}
  z_j &\sim \mathcal{N}(0, 1), \quad j = 1,\ldots,N, \\
  \log\{p_{Am}\}_j &= \mu_{\log p_{Am}} + \sigma_{\log p_{Am}} \cdot z_j.
\end{align}

This separates the population-level parameters ($\mu$, $\sigma$) from
individual-level deviations, greatly improving HMC sampling efficiency.

\subsection{DEBtox: Damage Dynamics}

The 4-state DEBtox model adds a scaled internal damage variable $D_w$:
\begin{equation}
\frac{\mathrm{d}D_w}{\mathrm{d}t} = k_d\bigl(\max(C_w - z_w,\, 0) - D_w\bigr).
\end{equation}

At steady state, $D_w = \max(C_w - z_w, 0)$, so the stress factor is
$s = b_w \cdot D_w$.  The $\text{EC}_{50}$ is computed analytically in
\code{generated quantities} as $z_w + 0.5/b_w$.


% ============================================================
\section{Practical Advice}
\label{sec:advice}

\subsection{Choosing \texttt{adapt\_delta}}

Start with \code{adapt\_delta = 0.8} (default).  If \fn{bdeb\_diagnose}
reports divergent transitions, increase to 0.9, 0.95, or even 0.99.
Higher values slow sampling but improve geometric fidelity.

\subsection{When to Use Hierarchical Models}

Use \code{type = "hierarchical"} when:
\begin{itemize}[nosep]
  \item More than 3--4 individuals are available.
  \item You want to estimate population-level distributions.
  \item Some individuals have sparse data (partial pooling helps).
  \item You need predictions for new, unobserved individuals.
\end{itemize}

\subsection{Prior Sensitivity}

Always compare results under at least two plausible prior specifications.
Parameters that remain prior-dominated (posterior $\approx$ prior) are
weakly identified and should be reported as such.

\subsection{Computational Tips}

\begin{itemize}[nosep]
  \item Stan model compilation is slow the first time ($\sim$30\,s) but
    cached for subsequent fits.
  \item Use \code{parallel\_chains = 4} to run chains in parallel.
  \item For large hierarchical models, start with a subset of individuals
    to debug parameterisation, then scale up.
  \item If ODE solver fails, check initial conditions (\code{E0}, \code{L0})
    and ensure they are biologically reasonable.
\end{itemize}


% ============================================================
\section{Comparison with Existing Tools}
\label{sec:comparison}

\begin{table}[htbp]
\centering
\caption{Comparison of \pkg{BayesianDEB} with existing DEB calibration tools.}
\label{tab:comparison}
\begin{tabular}{lcccc}
\toprule
\textbf{Feature} & \textbf{AmP/DEBtool} & \textbf{deBInfer} & \textbf{Custom Stan} & \textbf{BayesianDEB} \\
\midrule
Bayesian inference      & --          & Yes & Yes & Yes \\
Pre-built DEB models    & Yes         & --  & --  & Yes \\
Hierarchical models     & --          & --  & Manual & Built-in \\
DEBtox (TKTD)           & Partial     & --  & Manual & Built-in \\
Prior specification API & --          & Basic & Manual & Rich \\
Diagnostics             & Minimal     & Basic & Manual & Automated \\
PPC                     & --          & --  & Manual & Built-in \\
\R{} package            & MATLAB+\R{} & Yes & --  & Yes \\
\bottomrule
\end{tabular}
\end{table}


% ============================================================
\section*{Acknowledgements}

This package is based on the theoretical framework developed in the
\textit{Bayesian DEB (BDEB)} monograph.  The Stan Development Team
provided the computational infrastructure that makes modern Bayesian
DEB modelling practical.

% ============================================================
\begin{thebibliography}{9}

\bibitem[Carpenter et~al.(2017)]{Carpenter2017}
Carpenter, B., Gelman, A., Hoffman, M.D., et~al. (2017).
\newblock Stan: A probabilistic programming language.
\newblock \textit{Journal of Statistical Software}, 76(1), 1--32.

\bibitem[Gelman et~al.(2020)]{Gelman2020}
Gelman, A., Vehtari, A., Simpson, D., et~al. (2020).
\newblock Bayesian workflow.
\newblock \textit{arXiv preprint arXiv:2011.01808}.

\bibitem[Kooijman(2010)]{Kooijman2010}
Kooijman, S.A.L.M. (2010).
\newblock \textit{Dynamic Energy Budget Theory for Metabolic Organisation}.
\newblock Cambridge University Press, 3rd edition.

\end{thebibliography}

\end{document}
